\section{The computation being optimized}
\label{sec:C04:operator}
Attention kernels optimize a specific tensor operator; they do not determine what an agent should remember. For one head, let $Q\in\mathbb R^{n_q\times d_h}$, $K\in\mathbb R^{n_k\times d_h}$, and $V\in\mathbb R^{n_k\times d_v}$. With additive mask $A$, attention is $O=PV$, where $P=\operatorname{softmax}(QK^\top/\sqrt{d_h}+A)$ row by row. The mask can encode causality, padding, sliding windows, or a more specialized visibility relation. Changing the mask changes the operator. Changing its implementation while preserving the same mask and arithmetic target is an optimization of that operator. These cases require different correctness tests~\cite{vaswani2017attention,dao2022flashattention}.

For dense equal-length sequences and $d_v=d_h$, forming scores and multiplying probabilities by values requires approximately $4n^2d_h$ floating-point operations when one multiply-add counts as two operations. Causal structure can reduce useful score work by roughly half, but actual savings depend on whether the kernel skips masked tiles. Projection and feed-forward costs are additional. A dense kernel that uses less temporary memory still evaluates a quadratic set of interactions. FlashAttention is exact attention in this algorithmic sense; it avoids materializing the full score and probability matrices in high-bandwidth memory, rather than replacing them with a low-rank or sparse approximation~\cite{dao2022flashattention}.

Two regimes matter for agents. Prefill processes a supplied context and commonly has many query positions. Autoregressive decode often has one new query position per sequence, interacting with a long retained cache. A technique that improves a large square matrix computation need not produce the same gain for a thin decode operation. During a long tool pause, the cache may occupy accelerator memory without generating tokens. The workload consequently includes active computation, idle retained state, resumption, branching, and cancellation. Kernel benchmarks that omit those phases do not fully characterize a long-running agent service.

Correctness also has several levels. Algebraic equality assumes real arithmetic. Floating-point implementations reorder reductions and can differ numerically while targeting the same function. Randomized dropout requires consistent random-number semantics in training and recomputation. Mask alignment, positional offsets, and sequence lengths must match exactly. A kernel can pass average numerical tolerances yet fail at a rare boundary length or unsupported head dimension. For an agent, such a defect can appear as a semantic planning error many actions later. Operator-level verification should therefore precede task-level evaluation rather than relying on an aggregate task score to detect low-level mistakes.

